\documentclass[11pt,a4paper]{article}
\usepackage[margin=22mm]{geometry}
\usepackage{amsmath,amssymb}
\usepackage{fontspec}
\usepackage{polyglossia}
\setdefaultlanguage{russian}
\setmainfont{DejaVu Serif}
\usepackage[unicode,hidelinks]{hyperref}

\allowdisplaybreaks[1]
\raggedright
\setlength{\parskip}{0.35em}
\newcommand{\dd}{\mathop{}\!\mathrm{d}}
\newcommand{\avg}[1]{\langle #1\rangle}
\newcommand{\ket}[1]{\lvert #1\rangle}
\newcommand{\bra}[1]{\langle #1\rvert}
\newcommand{\braket}[2]{\langle #1\mid #2\rangle}
\DeclareMathOperator{\Var}{Var}

\newcommand{\solutiontitle}[4]{%
  \hypersetup{pdftitle={Квантовая механика 1. Решения ДЗ №#1}}%
  \begin{center}
    {\small Квантовая механика 1 \enspace\textperiodcentered\enspace Семестр 1\par}
    \vspace{1ex}
    {\LARGE\bfseries Решения домашнего задания №#1\par}
    \vspace{0.5ex}
    {\large #2\par}
  \end{center}
  \noindent{\small Условия приведены в файле «#3», #4.}\par\medskip
}

\begin{document}
\solutiontitle{4}{Радиальные распределения и представления}{ДЗ2026(1)\_04.pdf}{с.~1}

В одном измерении используем пару преобразований
\[
 a(p)=\frac1{\sqrt{2\pi\hbar}}\int e^{-ipx/\hbar}\Psi(x)\,\dd x,\qquad
 \Psi(x)=\frac1{\sqrt{2\pi\hbar}}\int e^{ipx/\hbar}a(p)\,\dd p .
\]
Эта нормировка сохраняет норму \(L^2\).

\section{Двумерное состояние}

Для \(\Psi=C\rho e^{-\rho^2/(2a^2)}\sin\varphi,\ a>0\)
интегрирование с мерой \(\rho\dd\rho\,\dd\varphi\) даёт
\[
 1=C^2\!\int_0^{2\pi}\sin^2\varphi\,\dd\varphi
       \int_0^\infty \rho^3e^{-\rho^2/a^2}\,\dd\rho
   =C^2\frac{\pi a^4}{2}.
\]
Берём \(C=\sqrt{2/\pi}/a^2\).
Радиальная плотность \(w(\rho)=2\rho^3e^{-\rho^2/a^2}/a^4\).
После замены \(u=\rho^2/a^2\) мера \(w(\rho)\dd\rho\) равна
\(u e^{-u}\dd u\), откуда
\[
 \avg{\rho^n}=a^n\Gamma(2+n/2),\quad
 \avg{\rho}=\frac{3\sqrt\pi}{4}a,\quad
 \avg{\rho^2}=2a^2,\quad
 \Var(\rho)=a^2\left(2-\frac{9\pi}{16}\right).
\]
У \(w\) логарифмическая производная \(3/\rho-2\rho/a^2\)
меняет знак с плюса на минус при \(\rho_{\rm mp}=a\sqrt{3/2}\).
Сдвиг \(\varphi\mapsto\varphi+\pi\) даёт пространственную инверсию,
сохраняет плотность и меняет знак \(\boldsymbol\rho\) и \(\Psi\).
Отсюда \(\avg{\boldsymbol\rho}=\mathbf0,\ \avg{P}=-1\).

\section{Трёхмерное состояние}

В \(\Psi=Cre^{-r/(2a)}\sin\theta\cos\varphi\)
угловая функция \(n_x=\sin\theta\cos\varphi\) удовлетворяет
\(\int n_x^2\dd\Omega=4\pi/3\). Это следует из равенства интегралов
для трёх осей и \(n_x^2+n_y^2+n_z^2=1\).
При \(a>0\) норма равна \(C^2(4\pi/3)\int_0^\infty r^4e^{-r/a}\dd r
=32\pi a^5 C^2\). Поэтому \(C=(32\pi a^5)^{-1/2}\), а радиальная
плотность \(w(r)=r^4 e^{-r/a}/(24a^5)\).
Интегрируя по частям или используя \(\Gamma(n+1)=n!\), получаем
\[
 \avg{r^n}=a^n\frac{(4+n)!}{4!}\quad(n=0,1,2,\ldots),\qquad
 \avg{r}=5a,\quad \avg{r^2}=30a^2,\quad \Var(r)=5a^2.
\]
Поскольку \(\Psi=Cx e^{-r/(2a)}\) нечётна, \(\avg{P}=-1\).
Кроме того \(P\mathbf pP=-\mathbf p\), откуда в собственном
состоянии чётности \(\avg{\mathbf p}=-\avg{\mathbf p}=\mathbf0\).
Угловую зависимость раскладываем по
\(\cos\varphi=(e^{i\varphi}+e^{-i\varphi})/2\).
Два слагаемых ортогональны по интегралу по \(\varphi\), имеют
одинаковые нормы и отвечают значениям \(+1\) и \(-1\)
оператора \(-i\partial_\varphi\) из условия. Вероятность каждого
равна \(1/2\). Для размерного момента значения равны \(\pm\hbar\).

\section{Координата в импульсном представлении}

Поскольку \(i\hbar\,\partial_p e^{-ipx/\hbar}=x e^{-ipx/\hbar}\),
\[
 \bra p\hat x\ket\Psi=
 \frac1{\sqrt{2\pi\hbar}}\int e^{-ipx/\hbar}x\Psi(x)\,\dd x
 =i\hbar\,\frac{\dd}{\dd p}a(p).
\]
Вывод применим на функциях Шварца, а оператор в \(L^2\)
получается переносом области определения \(\hat x\)
унитарным преобразованием. В трёх измерениях
\(\hat x_j=i\hbar\partial_{p_j}\).

\section{Коммутаторы координат и импульсов}

\textbf{(a)} Для целого \(n\ge1\) правило Лейбница даёт
\((xf)^{(n)}=xf^{(n)}+nf^{(n-1)}\). Поэтому в координатном
представлении
\[
 [p^n,x]f=(-i\hbar)^n n f^{(n-1)}=-i\hbar n p^{n-1}f,\qquad
 [x^n,p]f=i\hbar n x^{n-1}f.
\]
В импульсном представлении \(p\) умножает, а \(x=i\hbar\partial_p\).
Правила \((p^n f)'=p^nf'+np^{n-1}f\) и
\((pf)^{(n)}=pf^{(n)}+nf^{(n-1)}\) дают соответственно
\(-i\hbar np^{n-1}f\) и \(n(i\hbar)^n f^{(n-1)}
=i\hbar n x^{n-1}f\). При \(n=0\) оба коммутатора равны нулю.

\textbf{(b)} По правилу производной
\([p_i,x_k]=-i\hbar\delta_{ik}I\).
Для скалярного квадрата \(p^2=\sum_j p_j^2\) правило Лейбница даёт
\[
 [p^2,x_k]=\sum_j\bigl(p_j[p_j,x_k]+[p_j,x_k]p_j\bigr)
 =-2i\hbar p_k .
\]
Аналогично \([p_i,\mathbf r^2]=-2i\hbar x_i\).
Если \(x^2\) означает квадрат одной компоненты, общая формула
\([p_i,x_k^2]=-2i\hbar\delta_{ik}x_k\) даёт нужный частный случай.

\section{Преобразование заданных координатных функций}

Обозначим \(\gamma=\sqrt{2\pi\hbar}\).

\textbf{(a)} Пусть фиксированный импульс в синусе равен \(p_*\),
а \(p\) есть переменная преобразования. Разложение синуса на две
экспоненты и \(\int e^{-ipx/\hbar}\dd x=2\pi\hbar\delta(p)\) дают
\[
 a(p)=\frac{C\gamma}{2i}\bigl[\delta(p-p_*)-\delta(p+p_*)\bigr].
\]
При \(p_*\ne0\) среднее по периоду от \(\sin^2\) равно \(1/2\),
поэтому \(\|\Psi\|^2\) на \(\mathbb R\) расходится. На отрезке
длины \(L\), содержащем целое число периодов, усечённый синус
нормируется при \(|C_L|=\sqrt{2/L}\).
При \(p_*=0\) функция тождественно нулевая.

\textbf{(b)} Подстановка двух дельта-функций в интеграл даёт
\[
 a(p)=\frac{C}{\gamma}\bigl(\sqrt3\,e^{ipx_0/\hbar}
                                      -e^{-2ipx_0/\hbar}\bigr).
\]
Исходная функция есть распределение, для которого обычная норма
\(L^2\) не определена. Для примера заменим каждую дельта-функцию
прямоугольником ширины \(\varepsilon\) и площади один.
При \(x_0\ne0\) достаточно малые прямоугольники не пересекаются,
и их взвешенная сумма имеет квадрат нормы \(4/\varepsilon\),
так что \(|C_\varepsilon|=\sqrt\varepsilon/2\).
При \(x_0=0\) остаётся одна дельта-функция с весом \(\sqrt3-1\).

\section{Преобразование заданных импульсных функций}

\textbf{(a)} Обратное преобразование двух экспонент даёт
\[
 \Psi(x)=\frac{C\gamma}{2i}
           \bigl[\delta(x+2x_0)-\delta(x-2x_0)\bigr].
\]
При \(x_0\ne0\) синус в импульсном представлении имеет бесконечную
норму на всей оси. На интервале импульсов длины \(L_p\), кратной
периоду, усечённое состояние нормируется при \(|C|=\sqrt{2/L_p}\).
При \(x_0=0\) состояние нулевое.

\textbf{(b)} Из свойства дельта-функции
\[
 \Psi(x)=\frac{C}{\gamma}\bigl(e^{ip_0x/\hbar}-e^{2ip_0x/\hbar}\bigr).
\]
При \(p_0\ne0\) обычная нормировка невозможна. При замене дельта-функций
непересекающимися прямоугольниками площади один и ширины
\(\varepsilon\) по импульсу норма в квадрате равна
\(2|C_\varepsilon|^2/\varepsilon\), поэтому
\(|C_\varepsilon|=\sqrt{\varepsilon/2}\).
При \(p_0=0\) два слагаемых сокращаются.

\section{Прямоугольные пакеты по импульсу (7*)}

В обозначениях листка \(k\) и \(\delta k>0\) имеют размерность
импульса. По заданному обратному преобразованию
\[
 a_k(p)=\frac{\mathbf 1_{I_k}(p)}{\sqrt{2\delta k}},
 \qquad I_k=[k-\delta k,k+\delta k].
\]
Планшерель даёт норму один. При общей ширине для двух центров
\[
 \braket{\Psi_k}{\Psi_{k'}}=
 \frac{\lvert I_k\cap I_{k'}\rvert}{2\delta k}
 =\frac{\max(0,2\delta k-|k-k'|)}{2\delta k}.
\]
Ортогональность имеет место при \(|k-k'|\ge2\delta k\).
Для разных ширин в числителе также длина пересечения,
а знаменатель равен квадратному корню из произведения длин.
Условие \(\delta k\ll k\) само по себе пересечения не исключает.
Подстановка \(p=k+u\) сводит моменты к интегралам с равномерной
плотностью на \([-\delta k,\delta k]\).
\[
 \avg{p}=k,\qquad
 \avg{p^2}=k^2+\frac1{2\delta k}\int_{-\delta k}^{\delta k}u^2\,\dd u
 =k^2+\frac{(\delta k)^2}{3}.
\]

\end{document}
